\documentclass[11pt,a4paper]{article}
\usepackage[OT1]{fontenc}
\renewcommand{\rmdefault}{cmr}
\renewcommand{\sfdefault}{cmss}
\renewcommand{\ttdefault}{cmtt}
\usepackage{amsmath,amssymb,amsthm,mathtools}
\usepackage[margin=27mm]{geometry}
\usepackage{microtype,booktabs,longtable,array}
\usepackage[hidelinks]{hyperref}
\setlength{\emergencystretch}{2em}
\setlength{\parskip}{3pt plus 1pt}
\setcounter{tocdepth}{1}
\newtheorem{theorem}{Theorem}[section]
\newtheorem{lemma}[theorem]{Lemma}
\newtheorem{proposition}[theorem]{Proposition}
\newtheorem{corollary}[theorem]{Corollary}
\theoremstyle{definition}
\newtheorem{definition}[theorem]{Definition}
\newtheorem{example}[theorem]{Example}
\newtheorem{problem}[theorem]{Open problem}
\theoremstyle{remark}
\newtheorem{remark}[theorem]{Remark}
\newcommand{\N}{\mathbb N}
\newcommand{\NN}{\mathbb N_{>0}}
\newcommand{\Ftwo}{\mathbb F_2}
\newcommand{\lv}{\operatorname{lv}}
\newcommand{\norm}{\operatorname{N}}
\newcommand{\sz}{\operatorname{sz}}
\newcommand{\can}[1]{\mathopen{[\![}#1\mathclose{]\!]}}
\newcommand{\mle}{\preccurlyeq}
\newcommand{\mslt}{\prec}
\newcommand{\PhiF}{\Phi}
\newcommand{\Esrc}{\mathcal E_{\mathrm{src}}}
\newcommand{\headpart}{\operatorname{lim}}
\newcommand{\wt}{\operatorname{wt}}
\newcommand{\Ty}{\operatorname{Ty}}
\newcommand{\FV}{\operatorname{FV}}
\newcommand{\NF}{\operatorname{NF}}
\newcommand{\sgn}{\operatorname{sign}}
\newcommand{\id}{\operatorname{id}}
\newcommand{\BH}{\mathsf{BH}}
\newcommand{\ord}{\operatorname{ord}}
\newcommand{\bodyurl}[1]{\url{#1}}
\numberwithin{equation}{section}

\newcommand{\rk}{\rho}
\newcommand{\markof}[1]{#1^{\star}}
\newcommand{\ang}[2]{#1\langle #2\rangle}
\newcommand{\Acc}{\mathsf{Acc}}
\newcommand{\It}{\mathsf I}
\title{A Type-Confined Termination Ordering for G\"odel's T\\
Exact comparison, a ranked path-order restriction, and Bachmann--Howard height}
\author{}
\date{28 October 2026\\Working manuscript}
\begin{document}
\maketitle
\begin{abstract}
We construct an effective, context- and substitution-stable termination
ordering for the literal lambda formulation of G\"odel's T.  It is a restriction
of recursive StarHorpo by inclusion of the finite types occurring in a term.
The restriction is essential for the ordinal analysis.  A finite
target-observation interpretation decides the recursive comparison despite
the unbounded extra-argument lists generated by its Select clause.
Every original beta, case and recursion rule decreases strictly.  On a
disjoint finite ranked ground signature the restriction is exactly multiset
recursive path ordering.  Every individual descending rank is below the
Bachmann--Howard ordinal and their supremum is that ordinal.  The lower
argument builds literal native-recursion realizers for finite-type ordinal
iteration; the upper argument uses a finite-type arithmetic accessibility
proof and the published inductive-norm bound for $\mathrm{ID}_1$.
The ranked first-order restriction is distinguished from uncurrying the native
case functional, and descending height is distinguished from a linear order
type.  The construction does not decide comparison of the numerical
Wilken--Weiermann canonical assignments.
\end{abstract}
\tableofcontents

\section{A syntactic ordering question}
The normalization of G\"odel's T does not by itself give a convenient
syntactic ordering of its terms.  Wilken and Weiermann ask for extensions of
recursive path ordering adapted to the finite-type calculus and anticipate
Bachmann--Howard strength \cite[Introduction]{ww}.  A related question asks
whether comparisons between their particular ordinal-vector assignments can
be decided.  Those are different specifications.  The present paper concerns
the syntactic ordering specification.

There are three requirements whose proofs should not be conflated.  A
recursive description of a relation must actually give a terminating
comparison algorithm.  A well-founded relation must orient the original
typed rewrite rules and survive insertion in contexts and substitution.
Finally an ordinal-height theorem must concern that same relation, rather
than a better-behaved auxiliary ordering.  We establish these requirements
separately and then combine them.

\subsection{Terms and the first-order embedding}
Simple types are generated by a base type $o$ and arrows.  The type level is
$\lv(o)=0$ and $\lv(\sigma\to\tau)=\max(\lv(\sigma)+1,\lv(\tau))$.
The native constants are $0:o$, $S:o\to o$,
$D^\tau:o\to\tau\to\tau\to\tau$, and
$R^\tau:o\to(o\to\tau\to\tau)\to\tau\to\tau$.
Application is left associative.  The non-contextual rules are
\[
 (\lambda x.A)B\to A[x:=B],\quad
 D^\tau0AB\to A,\quad D^\tau(Sn)AB\to B,
\]
\[
 R^\tau0AB\to B,\qquad
 R^\tau(Sn)AB\to A\,n\,(R^\tau nAB).
\]
Substitution is capture avoiding.  The reduction relation includes their
instances in every well-typed term context.

All these constants are ranked nodes of arity zero, even when their simple
type is functional.  Encode an application $UV$ literally as
$@_{\sigma,\tau}(U,V)$, with declaration
$[\sigma\to\tau,\sigma]\to\tau$.  Denote this encoding by $\mu$.
It commutes with substitution by structural induction.  No normalization,
eta expansion or uncurrying is part of $\mu$.

Independently adjoin a finite ranked signature $\mathcal F$ with declarations
$[o,\ldots,o]\to o$.  These are actual ranked nodes; a ternary $d\in\mathcal F$
is not the curried native constant $D^o$.  Give all heads multiset status.
The typed application heads form one precedence-equivalence class, all native
constants form one strictly higher class, and the finitely many heads of
$\mathcal F$ have a fixed strict total precedence above those classes.

\begin{definition}[Type support]
$\Ty(s)$ is the set of literal subterm types and binder types of $s$, closed
under taking type subexpressions.  It is a finite subtype-closed set.
\end{definition}
The declaration of a ranked ground node contributes only its argument and
result types.  Its absent curried counterpart is not a subterm type.

\section{The recursive comparison}
We specialize the ordinary, bottom-free, no-meta-variable fragment of Kop's
recursive StarHorpo \cite[Definition 5.19]{kop}.  The details below fix the
fragment, so later comparisons do not hide an iterative reachability oracle.

An ordinary ranked head $f(\vec a)$ has its fixed base arguments.  A marked
head $f^\star(\vec a;\vec e):\tau$ retains them, may have a finite list of
additional arguments, and may be retyped to an output type $\tau$.
Retyping never changes an argument's type.  Marking an ordinary term passes
through its initial lambda spine and marks its first functional head.  If
there is no such head it is undefined.

\begin{definition}[Simultaneous special application]\label{def:angle}
For $q:\rho_1\to\cdots\to\rho_n\to\tau$, define
$\ang q{a_1,\ldots,a_n}$ as follows.  With no arguments the result is $q$.
If at least $n$ initial lambdas are present, consume exactly $n$ and substitute
their arguments simultaneously.  Otherwise consume the $k<n$ initial lambdas.
If the remaining \emph{original} head is functional, substitute in its existing
arguments, mark it if needed, and append all remaining $n-k$ arguments at once.
If that head is a variable, the operation is undefined.
\end{definition}
For example, with $T=o\to o$, $h:T$ a native head and $y:o$,
\[
 \ang{(\lambda x:T.x)}{h,y}\ \text{is undefined},\qquad
 \ang{\ang{(\lambda x:T.x)}h}y=h^\star(;y).
\]
Thus full special application is not obtained by iterating unary tables.

Write $\succeq_*$ for weak comparison of equal-typed terms and
$s\succ_*t$ for $s^\star\succeq_*t$.  The all-Mul ordinary clauses are:
\begin{enumerate}
\item $x\succeq_*x$.
\item $\lambda x.s\succeq_*\lambda x.t$ if $s\succeq_*t$.
\item $f(\vec s)\succeq_*g(\vec t)$ if $f\approx g$ and their equally long
argument multisets admit a bijection of weak comparisons.
\item $f(\vec s)\succeq_*t$ if $f^\star(\vec s)\succeq_*t$.
\item A marked source $P=f^\star(\vec s;\vec e):\tau$ weakly dominates $t$
if an argument $q$ from either list has a suffix type
$\rho_1\to\cdots\to\rho_n\to\tau$ and
$\ang q{P_{\rho_1},\ldots,P_{\rho_n}}\succeq_*t$.
The case $n=0$ is included.
\item $P\succeq_*\lambda x.t$ if, with fresh $x$, the marked source with $x$
appended to its extras weakly dominates $t$.
\item $P\succeq_*g(\vec t)$ if $f$ is strictly above $g$ and the source
retyped to each target-argument type weakly dominates that argument.
\item The same conclusion holds when $f\approx g$, the \emph{base} argument
multiset of $P$ strictly exceeds the target multiset under the generalized
weak/strict multiset test, and those whole-source premises also hold.
\end{enumerate}
These are respectively Var, Abs, Fun, Put, Select, F-Abs, Copy and Stat.
All derivations are finite.  Strict multiset comparison retains weakly matched
occurrences and replaces or deletes a nonempty source part, with every new
occurrence strictly below an occurrence of that part.  Extras are selectable
but do not enter the Stat base multiset.

The published strong-pair theorem for this recursive relation supplies
transitivity, compatibility, well-foundedness and substitution stability
under a well-founded precedence and compatible statuses
\cite[Theorems 5.32, 5.37 and 5.39]{kop}.  We use these published properties;
they do not supply the new finite comparator or the calibration proved here.

\section{An exact finite comparison algorithm}\label{sec:comparator}
Fix ordinary inputs $(s,t)$.  Let $K$ be their subtype closure, $V$ their free
names, and $B$ the maximum nesting depth of a lambda in $t$.  Recursive target
premises are opened subterms of $t$.  A target lambda is opened at a fresh
typed port.  Use $B$ port slots; at context $\Gamma$ retain only opened
observations with remaining lambda depth at most $B-|\Gamma|$.
Write $O_\sigma^\Gamma$ for this finite set of observations of type $\sigma$.
Contexts extend by adding previously unused slots, keeping their types.

\begin{lemma}[Finite closure]\label{lem:finiteclosure}
Every type requested by a recursive comparison belongs to $K$, every source
head is a head of the original source, and every target observation and
opening context belongs to the finite family just described.
\end{lemma}
\begin{proof}
Select uses a suffix of the type of an existing argument; all its argument
and result types are subtypes of that type.  The retyped copies introduce
only these types.  Copy and Stat request target-argument types.  F-Abs uses
the domain and codomain of the observed target lambda.  Substitution does
not change a body-node type.  These exhaust the ways to create a type.
Marking and Select preserve the source head; Copy creates no new recursive
source with its target head.  The only fresh port is introduced while
descending a target lambda, so used ports plus remaining target depth never
exceeds $B$.  The free names of generated sources lie in $V\cup\Gamma$.
\end{proof}

We now give finite domains on which special application and source comparison
can be computed without bounding the size of generated sources.

\begin{definition}[Natural finite domains]\label{def:domains}
At base type, a value at $\Gamma$ consists, at each extension $\Delta$, of
an availability bit and weak/strict subsets of $O_o^\Delta$.
At arrow type $\sigma$, add for each nonempty suffix decomposition
$\sigma=\rho_1\to\cdots\to\rho_n\to\tau$ a monotone full tuple table
\[
 D_{\rho_1}^\Delta\times\cdots\times D_{\rho_n}^\Delta
       \longrightarrow D_\tau^\Delta.
\]
All argument and result types are proper subtypes, so this is a recursion
on types.  Values are ordered pointwise.  Weakening $a\mapsto a|_\Theta$
retains the components at extensions of $\Theta$.
Restrict the raw finite products to values satisfying every equation
\[
 \operatorname{App}_\Delta(a,\vec v)|_\Theta
   =\operatorname{App}_\Theta(a|_\Theta,\vec v|_\Theta)
       \qquad(\Delta\subseteq\Theta).
\]
Only these hereditary natural values form $D_\sigma^\Gamma$.
\end{definition}

\begin{lemma}[Effective complete lattices]\label{lem:lattices}
The domains in \ref{def:domains} are finite complete lattices, effectively
enumerable from the inputs.  Weakening preserves their joins and meets.
\end{lemma}
\begin{proof}
The claim holds for finite products of bits and powersets.  At an arrow,
the lower-type domains are already constructed.  Enumerate all tables and
retain the monotone ones satisfying the finitely many naturality equations.
Pointwise joins and meets of a family of such tables remain monotone and
natural, because weakening is a coordinate projection and preserves lower-type
joins and meets.  The bottom and top families satisfy the equations as well.
This proves the inductive claim and gives a finite construction.
\end{proof}

For every possible source head $f$, output type $\tau\in K$, and context,
introduce a monotone constructor table $M_{f,\tau}^\Gamma(\vec b,E)$.
The tuple $\vec b$ has the fixed base types of $f$; $E$ is a down-set in
the disjoint union of the finitely many typed domains at $\Gamma$.
Require constructor coherence
\begin{equation}\label{eq:constructorcoherence}
 M_{f,\tau}^\Gamma(\vec b,E)|_\Delta
 =M_{f,\tau}^\Delta(\vec b|_\Delta,
          {\downarrow}\{a|_\Delta:a\in E\}).
\end{equation}
The coherent monotone families form a finite complete lattice, by the same
pointwise argument.  The down-set represents the values of \emph{all} extra
arguments, without preserving their irrelevant order or multiplicity.

At the syntactic level, irrelevance of extra order and repetition can be
checked directly. Let $\equiv_E$ be congruence generated by permuting or
repeating marked extras, keeping every distinct extra and every base
argument. It preserves types, free names, substitution, retagging and the
original-spine operation. Induction on a finite comparison derivation
transports it across $\equiv_E$: Select still finds its corresponding
argument, F-Abs appends the same fresh variable, and Copy and Stat inspect
only the unchanged head and base list. Var, Abs, Fun and Put use their
induction premises. The reverse implication follows by symmetry.
This explains the set convention but does not prove semantic compression
by itself; adequacy below proves that the finite down-set abstraction also
loses no comparison.

\subsection{The fixed-point equations}
For a marked state $(f,\vec b,E)$ use the following equations for its weak
observation bits.  All data are first weakened to the observation context.
\begin{enumerate}
\item Select chooses a value $a$ from $\vec b$ or $E$ and a compatible
suffix decomposition.  Apply its one full tuple table to the corresponding
retyped $M$ values of the whole state; use the resulting weak bit.
\item Copy tests strict head precedence and the weak bits for every
target argument of the correspondingly retyped whole state.
\item Stat tests equivalent heads, the finite generalized strict
multiset test on $\vec b$, and the same whole-source bits.
\item F-Abs opens the target at a fresh port, weakens existing inputs,
and appends the port value to the extra down-set.
\end{enumerate}
The availability bit is true.  A marked state's strict bits are empty.
Its full tuple table appends the entire supplied tuple to $E$ and consults
the result-type constructor.  An unmarked $f(\vec b)$ has weak bits from
Fun or its marked state, strict bits from that marked state's weak bits,
and tuple tables that mark the head and append all supplied arguments.
Variables have their literal weak bits, no strict bits, and undefined
nonempty tuple tables.  Undefined means the whole bottom value, including
availability, not just an empty weak observation set.

A lambda has matching bits against lambdas obtained from its opened body.
For its tuple table, scan the \emph{original} lambda spine.  Each consumed
binder is recorded as protected.  With arguments left at a protected-variable
head return bottom; at a functional head append all remaining arguments
simultaneously; at an unprotected free name consult its full table.
With no arguments left interpret the residual term in the updated environment.
This is a structural helper, not extensional beta normalization.

These equations define an operator $F$ on the finite constructor lattice.
\begin{lemma}[Structural interpretation]\label{lem:structural}
At a natural coherent constructor family, interpretation and the original-spine
helper are monotone, commute with weakening and capture-avoiding substitution,
and interpret simultaneous special application exactly.
\end{lemma}
\begin{proof}
Induct jointly on the finite source and the original-spine helper.
At a variable use its assigned natural value; the protected variable case is
constantly bottom.  At a functional head use monotonicity and
\eqref{eq:constructorcoherence}.  At a lambda consume a binder with a fresh
name outside $V$, the port supply and substitution ranges, and apply the
induction hypothesis to the original body.  Its literal lambda matching bits
are retained independently of its tuple behavior.  A free substituted variable
uses its full tuple table, which has already recorded its own original spine.
If the supplied tuple ends within a lambda spine the two constructions are
ordinary simultaneous substitution; otherwise the original residual head
determines either the same marked head with the same extras or the same
undefined case on both sides.  Naturality at the variable case is precisely
the tuple equation in \ref{def:domains}; at a functional head it is
\eqref{eq:constructorcoherence}.  No unary-currying identity is used.
\end{proof}

\begin{lemma}[A finite monotone operator]\label{lem:operator}
$F$ preserves natural coherent families and is monotone.  Iteration from
bottom reaches its least fixed point in finitely many steps.
\end{lemma}
\begin{proof}
Every rule test is positive in weak and strict bits.  Application uses monotone
tables, weakening is a projection, and taking a down-set is monotone.
The structural interpretation is monotone by \ref{lem:structural}.
At extension $\Delta$ every observation equation consults only weakened
inputs and $M^\Delta$.  For tuple components use
\[
 {\downarrow}(E\cup\{v_i\})|_\Delta
 ={\downarrow}(E|_\Delta\cup\{v_i|_\Delta\}),
\]
which follows from monotonicity of weakening and inclusion of the original
image on each side.  This proves coherence and tuple naturality of $F(M)$.
The ascending bottom iteration stabilizes because the lattice is finite.
\end{proof}

\begin{lemma}[Adequacy]\label{lem:adequacy}
At the fixed point, a weak or strict observation bit is true exactly when the
corresponding finite recursive comparison is derivable.
\end{lemma}
\begin{proof}
For soundness relate a value $a$ at $\Gamma$ to a concrete term $q$ with
$\FV(q)\subseteq V\cup\Gamma$.  Require that all its bits at all extensions
are derivable comparisons from $q$, and that each full tuple table preserves
the relation on related, defined, concrete argument tuples.  Only bottom
relates to an undefined result.  This relation is downward closed, by induction
on types.  Bottom relates to every permitted defined term.

Prove by stage induction that each constructor approximation preserves this
relation.  Select picks a concrete base or extra argument related to the
selected value; a value below an actual extra's value remains related to that
same extra.  The tuple property and \ref{lem:structural} give its literal
premise, so the Select rule gives the conclusion.  Copy and Stat use their
displayed premises.  F-Abs uses a fresh port, fresh for the source by the
support condition; its enlarged context supplies precisely the opened
premise.  Tuple outputs are the corresponding concrete marked head with all
supplied arguments appended.  The structural fundamental lemma is proved by
the same source/helper induction as \ref{lem:structural}.  Thus the limiting
interpretation is sound.

For completeness induct on a finite comparison derivation.  Var, Abs and Fun
are their structural equations.  Put inserts a marked weak bit into an
unmarked weak set, and strictness inserts that bit into its strict set.
Select uses the exact full-tuple operation; every actual extra is in its
downward closure, so it is not lost.  Copy, Stat and F-Abs are their equations.
The context/type invariants in \ref{lem:finiteclosure} retain every target
premise.  These exhaust the bottom-free, no-meta-variable rules.
\end{proof}

\begin{theorem}[Effective comparison]\label{thm:decide}
Both recursive comparisons on ordinary bottom-free inputs are decidable.
\end{theorem}
\begin{proof}
Construct the finite domains and constructor lattice, compute its fixed point,
interpret $s$ at the empty context, and test the root observation $t$.
Adequacy proves both answers correct.  No generated-source size bound or
numerical Select fuel is part of the algorithm.
\end{proof}

\section{Confinement, stability and exact ranked RPO}
\begin{definition}
For equal outer types put
\[
 s\succeq_c t\iff\mu(s)\succeq_*\mu(t)\ \land\ \Ty(t)\subseteq\Ty(s),
 \qquad
 s\succ_c t\iff\mu(s)\succ_*\mu(t)\ \land\ \Ty(t)\subseteq\Ty(s).
\]
\end{definition}
\begin{lemma}[Free-variable containment]\label{lem:fv}
Every weak comparison $s\succeq_*t$ satisfies $\FV(t)\subseteq\FV(s)$.
\end{lemma}
\begin{proof}
Use simultaneous induction for weak comparison and strictness through marking.
Var is immediate; Abs removes the common opening variable.  Fun distributes
target arguments to source arguments.  Put preserves free variables.
Copy and Stat have a whole-source premise for each target argument.
In Select the special application has free variables only from its selected
argument and the supplied whole-source copies, all already in the source.
F-Abs adds a fresh opening variable to the extras, and binding it in the
target removes exactly that possible addition.  No remaining rule applies.
\end{proof}
\begin{lemma}[Exact support substitution]\label{lem:supportsub}
For a finite type-preserving substitution $\gamma$,
\[
 \Ty(s\gamma)=\Ty(s)\cup\bigcup_{x\in\FV(s)}\Ty(\gamma(x)),
\]
where absent entries mean the variable itself.
\end{lemma}
\begin{proof}
Every old node retains its result or binder type, and every new node is in an
inserted image.  Conversely every type in an inserted image occurs at an
actual free occurrence.  Structural induction gives both inclusions.
\end{proof}
\begin{theorem}[Stable effective pair]\label{thm:stable}
The confined pair is decidable, compatible and transitive; its weak part is
reflexive and its strict part is well-founded.  Both are stable under
well-typed substitution and every literal context.
\end{theorem}
\begin{proof}
Decidability follows from \ref{thm:decide} and a finite support test.
The published StarHorpo strong-pair theorem gives the underlying relational
properties.  Support inclusions compose, so restriction preserves
transitivity and both mixed compatibilities.  Restriction preserves
well-foundedness.  A context adds the same set of types to each side.
For substitution use \ref{lem:fv} and \ref{lem:supportsub}: every image type added
on the target side is also added on the source side.  Combine this with the
published substitution property and $\mu(s\gamma)=\mu(s)\gamma_\mu$.
\end{proof}

\begin{theorem}[Exact ranked restriction]\label{thm:rpo}
On the first-order $\mathcal F$-algebra, $\succ_c$ is precisely multiset RPO
for the selected precedence.  Its weak part is RPO together with recursive
multiset equality.
\end{theorem}
\begin{proof}
Every type support is $\{o\}$.  Between first-order endpoints, an unmarked
source uses only Var, Fun or Put; an initially marked source uses only Select,
Copy or Stat.  Every argument is ground, so Select has zero supplied
arguments.  No extras or lambdas are generated.  The recursive comparisons
therefore stay within this first-order fragment.

Induct on a weak derivation, interpreting an unmarked comparison as weak RPO
and an initially marked comparison as strict RPO.  Var is equality.  Fun is
a permutation of weak argument comparisons: either all are recursively equal,
or the multiset strictly decreases.  Put includes strictness into weakness.
Select is the subterm clause, Copy is the precedence clause and Stat is the
strict multiset clause.  Conversely induct on an RPO/equality derivation.
Use Var/Fun for recursive multiset equality, zero-argument Select for the
subterm clause, Copy for precedence, and Stat for multiset decrease.
Stat's side comparisons with target arguments follow by Select from the
matched source arguments, using Put to include any strict match weakly.
There is no derivation path through a hidden higher-order intermediate.
\end{proof}

\begin{example}[A ranked comparison]
Take a ranked constant $a:o$, heads $f:[o,o]\to o$ and $g:[o]\to o$,
and precedence $f>g>a$. Then
$f(a,g(a))\succ_c g(g(a))$. The marked $f$ source selects its second
argument with zero supplied arguments, giving the strict subterm
comparison with $g(a)$. Copy at $f>g$ uses that whole-source premise
to construct the target $g(g(a))$. All type supports are $\{o\}$, so
confinement imposes no additional test. The theorem identifies precisely
the same derivation with ordinary multiset RPO; no curried application
node is involved.
\end{example}

\section{The original reductions decrease}
\begin{lemma}[Template grafting]\label{lem:graft}
For an ordinary template $A$ and $p\succeq_*q$ of the same type,
$A[x:=p]\succeq_*A[x:=q]$, even if $p$ contains marked subterms.
\end{lemma}
\begin{proof}
Induct on the ordinary template.  Its $x$ leaf is the hypothesis; other
variables are Var.  Abs uses fresh bound names, and Fun uses the induction
comparisons with identity argument matching.  This is not a claim about
arbitrary marked contexts.
\end{proof}
\begin{lemma}[Beta and cases]\label{lem:betacases}
Every original beta and case step decreases under $\succ_c$.
\end{lemma}
\begin{proof}
For a beta redex put $P_\rho=@^\star(\lambda x.A,B):\rho$.
Select on its second base argument gives $P_\sigma\succeq_*B$.
Template grafting gives $A[x:=P_\sigma]\succeq_*A[x:=B]$.
Select on the first argument with its one self-copy now gives
$P_\tau\succeq_*A[x:=B]$, hence strict beta orientation.

For $D^\tau uAB$, select the outer second argument to obtain $B$.
For $A$, select the outer first argument with one self-copy; the resulting
marked inner application still has $A$ as its second base argument, selected
with zero arguments.  These comparisons hold for every test $u$, and in
particular orient the two original rules.  The right-hand-side type support
is included in that of the original redex in each case.
\end{proof}

\begin{lemma}[Explicit application contexts]\label{lem:context}
If ordinary $a\succ_*b$, replacing $a$ by $b$ in any ordinary context gives
another strict comparison.
\end{lemma}
\begin{proof}
Strict comparison implies weak comparison: peel common initial lambdas and
apply Put at the marked head.  At a ranked context head, Stat matches all
unchanged arguments reflexively and the changed argument strictly.  Each
whole-source side premise follows by zero-argument Select of the corresponding
source argument and its reflexive or weak comparison.  Abs handles a lambda
context.  Induct on context depth.  The support condition is preserved by
union, so the result also holds for $\succ_c$.
\end{proof}

For the successor recursor write $H=o\to\tau\to\tau$,
$U_1=R^\tau(Sn)$, $U_2=U_1A$, $L=U_2B$, and $L'=R^\tau nAB$.
Use the following marked families, with their original heads and base lists:
\[
 P_\rho=@^\star(U_2,B):\rho,\qquad
 C_{2,\rho}=@^\star(U_1,A;P_\tau):\rho,
\]
\[
 C_{1,\rho}=@^\star(R^\tau,Sn;C_{2,H},C_{2,\tau}):\rho,
 \quad C_{0,\rho}=(R^\tau)^\star(;C_{1,o},C_{1,H},C_{1,\tau}):\rho.
\]
\begin{lemma}[Recursor orientation]\label{lem:recursor}
Both original recursor rules decrease strictly under $\succ_c$.
\end{lemma}
\begin{proof}
The zero rule selects the outer second argument $B$.
For the successor rule, three Select prefixes carry a weak comparison from
$C_{0,\tau}$ back through $C_{1,\tau},C_{2,\tau}$ to $P_\tau$.
They select $U_2$, then $U_1$, then $R^\tau$, with all self-arguments supplied
simultaneously at the displayed types.

From $C_{0,H}$ select successively $C_{1,H}$, $C_{2,H}$ and $A$ to obtain
$C_{0,H}\succeq_*A$.  From $C_{0,o}$ select $C_{1,o}$ and $Sn$; since
$Sn\succ_*n$ by outer Select, this yields $C_{0,o}\succeq_*n$.
At type $\tau$, the extra chain reaches $P_\tau$.  Context closure of
$Sn\succ_*n$ gives $L\succ_*L'$, or $P_\tau\succeq_*L'$; prefixing
the extra chain gives $C_{0,\tau}\succeq_*L'$.
Two Copy clauses at $R^\tau>@$ first give
$C_{0,\tau\to\tau}\succeq_*An$, then
$C_{0,\tau}\succeq_*AnL'$.  Prefix the three outer Selects.
All types in the contractum already occur in the redex; this proves confinement.
The reasoning is uniform in arrow or base $\tau$.
\end{proof}
\begin{corollary}
Every original T reduction, in every well-typed context and substitution
instance, strictly decreases under $\succ_c$.
\end{corollary}

\section{An individual Bachmann--Howard upper bound}\label{sec:upper}
For a well-founded relation write
$\rk_R(x)=\sup_{xRy}(\rk_R(y)+1)$.
The ordinal-analysis input is the following published theorem: for a positive
arithmetic inductive operator, every individual membership atom provable in
$\mathrm{ID}_1$ has inductive norm below $\BH$; for the accessibility operator
this norm equals the descending rank
\cite[Lemma 6.5.2, Definition 9.1.2 and Theorems 9.5.5--9.5.6]{pohlers}.
We give the finite-type accessibility and edge interpretation needed to use it.

Fix a finite subtype-closed set $S$.  Number the marked heads by
$k\in\N\cup\{\omega\}$, where unmarked heads have label $\omega$.
Retain finite base and extra lists and output annotations in $S$, and add
$\bot_\sigma$ for each $\sigma\in S$ and ordinary auxiliary application.
Use the all-Mul numbered reduction of \cite[Section 5.3]{kop}, restricted to
$S$: marking $\omega$ to any finite label; Select at $k+1$ applying a chosen
argument to retyped self-copies at $k$; Copy to a lower-precedence unmarked
head with self-arguments at $k$; Mul replacement of a nonempty part of the
base list by marked arguments; creation of an abstraction with one extra and
label decreased by one; reductions in ranked arguments and lambda bodies;
function-side reduction in auxiliary application; beta; and reduction to
bottom.  Every instance is well typed in $S$.  There is no argument-side
monotonicity for auxiliary application.

The primary equivalence permutes equivalent heads and their appropriate
base arguments, preserving labels, types and extra-list order.  With de Bruijn
binders its classes are finite and computable: the shape is unchanged and only
finitely many $S$-heads and permutations occur.  Work on the least-code
representatives.  Let $H$ be the induced one-step relation and $R=H^+$.
Both are arithmetic; equivalence-only motion has been quotiented out.

\begin{lemma}[Arithmetic accessibility proof]\label{lem:id1}
For this fixed $S$, $\mathrm{ID}_1$ proves that every well-typed auxiliary
term is accessible under $R$.
\end{lemma}
\begin{proof}
Use two positive arithmetic operators: $I(x)$ says every $R$ successor of
$x$ is in $I$, and $J(M)$ says every direct multiset replacement successor
of the finite multiset $M$ is in $J$.  They are separate unnested inductive
definitions, with well-typedness included in both defining formulas.
They refer to the arithmetic relations $R$ and $\operatorname{DM}(R)$,
not to one another's inductive predicates. Disjoint numerical tags combine
them into one positive arithmetic definition, as permitted in
$\mathrm{ID}_1$. For the finitely many types in $S$, define arithmetic
computability predicates by $\mathsf C_o(t)=I(t)$ and
\[
 \mathsf C_{\sigma\to\tau}(t)\iff
       \forall u(\mathsf C_\sigma(u)\to\mathsf C_\tau(tu)).
\]
The type recursion is finite; no new inductive definition at an arrow is used.
Typedness is conjoined in each clause. Quantification over arbitrary
computable arguments in the following induction invariants abbreviates
the finite disjunction of these typed formulas over $S$; it is an actual
first-order formula, rather than a satisfaction predicate for all simple types.

First prove by type induction that a non-lambda is computable when all its
direct reducts are; variables and bottoms are computable; computability implies
$I$ and is closed under $H$ and $R$.  At base type use closure and transitivity
of $R$.  At an arrow, a non-lambda application $tu$ reduces only in its function
or to bottom, so the lower-codomain induction applies.  To infer $I(t)$ from
$I(tx)$ for a fresh variable, use $I$-induction: every descendant $t'$ of $t$
has $t'x$ as a descendant of $tx$.  This proves inverse application.

For a lambda whose computable instances $v[x:=u]$ are all computable, choosing
$u=x$ gives $I(v)$.  Use $I$-induction on $v$, retaining the universal instance
premise.  The reducts of $(\lambda x.v)u$ are its beta contractum,
function-side body reducts and bottom; all are computable by the premise,
the induction hypothesis and the bottom case.  Hence the application and
the lambda are computable.

Next, accessibility of the members implies $J(M)$.  To prove closure under
multiset union, use outer $J$-induction on one side and inner $J$-induction
on the other.  Label removed occurrences by their origin and assign each
replacement to a removed dominator; each component is unchanged or decreases,
and one decreases.  Then $I$-induction on $a$ gives $J(\{a\})$, since every
replacement is a finite collection of descendants of $a$.  Finite unions
give the general claim.

Prove that a numbered constructor with computable arguments is computable
by the lexicographic induction
\[
 (\text{head class},\text{base multiset},k,\text{extra multiset}).
\]
This means finite external induction on head classes, $J$-induction on the
base multiset with all later parameters universally quantified, induction
on finite $k$ followed by the $\omega$ case, and $J$-induction on extras.
Marking decreases $k$.  Select applies a computable selected argument to the
lower-label computable self-copies.  Copy uses lower $k$ then a lower head
class.  Mul decreases the base multiset because an initially marked
argument is an $R$ descendant of its unmarked argument; its new head may
reset to $\omega$ because the base parameter is earlier.  Abstraction creation
decreases $k$ before adding any extra, and uses the lambda fact just proved.
Argument reductions decrease the appropriate multiset.  Bottom is computable.
These are all auxiliary rule cases.  Finally structural induction under a
computable substitution makes every $S$-term computable, hence in $I$.
All induction invariants are first-order formulas of $\mathrm{ID}_1$.
\end{proof}

\begin{lemma}[Arithmetic simulation]\label{lem:simulation}
There is a computable translation $\operatorname{tr}$ such that
$s\succ_ct$ implies $\operatorname{tr}(s)R\operatorname{tr}(t)$ in the
fixed universe $S=\Ty(s)$.
\end{lemma}
\begin{proof}
Relate a recursive source $q$ to a numerical source $Q$ at budget $n$ when
unmarked heads correspond to label $\omega$ and every marked head corresponds
to a finite label at least $n$, with identical type, base and extra structure.
At quotient classes require a witnessing raw representative.  The relation
is arithmetic, decreases with $n$, and commutes with substitution.

If $q$ and $a_i$ correspond at budget $n+l$, where $l$ arguments are supplied,
the auxiliary application $QA_1\cdots A_l$ reduces to a source corresponding
at budget $n$ to the defined $\ang q{a_1,\ldots,a_l}$.  Consume the same original
lambda spine by beta; at its functional residual head mark $\omega$ if needed,
then create an abstraction and beta-contract once for each remaining argument.
Each such operation appends that argument and spends one label.  Definedness
was tested on the original residual head, before substituted binders are
inspected.  All types are domain or suffix types of $q$.

Induct on a finite recursive derivation.  Var needs budget zero; Abs and Fun
use the maximum child budget and context or quotient permutation.  Put marks
$\omega$ at a sufficiently large finite label.  Select with $l$ arguments uses
child budget plus $l+1$: its one Select step makes ordinary applications to
unchanged lower-label self-copies, and the preceding angle simulation supplies
the exact premise.  F-Abs uses child budget plus one.  Copy uses one plus the
largest child budget.  Stat chooses its finite retained/changed occurrence
assignment; changed source arguments are initially marked at the preceding
label and retained ones are unchanged, then child certificates finish every
target position.  Its budget is one plus their maximum.  All choices and
paths are finite arithmetic certificates.  An original strict comparison
starts by a nonempty marking step, so its complete simulated path is positive.
Every requested type is in $\Ty(s)$ by confinement and \ref{lem:finiteclosure}.
\end{proof}

\begin{theorem}[Individual upper bound]\label{thm:upper}
Every descending $\succ_c$ rank is strictly below $\BH$.
\end{theorem}
\begin{proof}
For a fixed source use $S=\Ty(s)$.  Every confined descendant stays in $S$.
The two preceding lemmas give an $\mathrm{ID}_1$ proof of accessibility of
$\operatorname{tr}(s)$ and a rank-monotone edge simulation.  The published
inductive-norm extraction gives
$\rk_{\succ_c}(s)\le\rk_R(\operatorname{tr}(s))<\BH$.
Only an individual membership atom is needed; no uniform notation for all
ranks or numerical bound on comparison time is inferred.
\end{proof}

\section{Native finite-type rank iteration}\label{sec:lower}
Restrict $\succ_c$ to closed terms of each type and write $\rk$ for its rank.
Restriction cannot increase a rank.  Its type-stable closed contexts retain
the properties already proved.  Put $\sigma=o\to o$.
All these ranks are countable because the syntax is countable and the
relation is well founded; no upper calibration is needed for that fact.

\begin{lemma}[A native carrier]\label{lem:carrier}
Let $P:B\to B$, $U:B$, and $K\succeq_cP$ be closed.  With
$W=R^B0(\lambda n:o.K)U$, one has $W\succ_cU$, and
\[
 W\succ_cT\quad\Longrightarrow\quad W\succ_cPT.
\]
\end{lemma}
\begin{proof}
Use the four marked families from the recursor proof with test $0$, step
$\lambda n.K$, and base $U$.  Their three Select prefixes lift weak
comparisons from the marked native $R^B$ to the marked outer $W$.
The extra chain reaches $U$.  At output $B\to B$, select the step-function
extra, supplying a ground self-copy; follow its original marked application
extras to $\lambda n.K$ and consume its unused ground binder.  The result is
$K\succeq_cP$.  At output $B$, the extra chain reaches $W^\star$, to which
the assumed $W^\star\succeq_*T$ is appended.  Copy at $R^B>@$ now constructs
$PT$, and prefixing the three Select steps gives the assertion.  Every target
type already occurs in $P,U,K,W$; the assumed confinement is preserved.
\end{proof}

Translate semantic ordinal types by $|\iota|=\sigma$ and arrows
homomorphically.  At the semantic base use the countable ordinals, with all
set functions at arrows.  Say $t\Vdash_{\iota}\alpha$ if $\rk(t)\ge\alpha$,
and $t\Vdash_{A\to D}f$ if every closed realizer $u$ of $x$ satisfies
$tu\Vdash_D f(x)$.  Countable suprema at arrows are pointwise.

\begin{lemma}[Realization closure]\label{lem:realclosure}
Realization is backward closed under $\succ_c$, downward closed in the
semantic pointwise order, and closed under countable suprema realized by
one fixed term.  The typed lambda fundamental lemma holds.
\end{lemma}
\begin{proof}
Induct on semantic type.  At base use the rank inequality and closure of
ordinal upper bounds under countable suprema.  At arrows apply the induction
claim to every realizing argument and use context closure.  For the lambda
case substitute the given realizers and use the oriented beta step.  Variables
and application are exactly the environment and arrow definitions.
\end{proof}

Let $E_{\iota}(i)=i$ and
$E_{A\to D}(i)=\lambda x:|A|.E_D(i)$.
For $B=|A|$, $\eta=B\to B$, and $Q_\eta=\lambda a:\eta.\lambda b:\eta.a$, put
\[
 K_i=Q_\eta P E_{A\to A}(i),\quad
 W_i=R^B0(\lambda n:o.K_i)U.
\]
\begin{lemma}[Iteration inequalities]\label{lem:iterationedges}
$W_i\succ_cU$.  If $i\succ_cj$ then
$W_i\succ_cW_j$ and $W_i\succ_cPW_j$.
\end{lemma}
\begin{proof}
The nonempty beta path gives $K_i\succeq_cP$.  The carrier lemma yields
$W_i\succ_cU$.  Literal context closure gives $W_i\succ_cW_j$; apply
the carrier lemma to obtain the last inequality.  The fixed template adds
the same types on both sides, and $PW_j$ adds only the already present
$B\to B$ and $B$.
\end{proof}

\begin{theorem}[Full-prefix ordinal iteration]\label{thm:iteration}
There is an effective closed native term $\It_A$ realizing
\[
 (\alpha,f,x)\longmapsto f^\alpha(x),\qquad
 f^0(x)=x,\quad f^{\beta+1}(x)=f(f^\beta(x)),\quad
 f^\lambda(x)=\sup_{\xi<\lambda}f^\xi(x).
\]
No monotonicity of $f$ is required.
\end{theorem}
\begin{proof}
Take $\It_A=\lambda i:\sigma.\lambda p:\eta.\lambda u:B.
R^B0(\lambda n:o.Q_\eta p E_{A\to A}(i))u$.
For closed realizers $P,U$, prove that the single $W_i$ realizes every orbit
value $x_\beta$ for $\beta\le\rk(i)$.  Induct externally on $\rk(i)$ and
internally on $\beta$.  Zero follows from $W_i\succ_cU$.
For $\gamma+1\le\rk(i)$, the rank recursion supplies $i\succ_cj$ with
$\rk(j)\ge\gamma$.  The outer induction gives $W_j\Vdash x_\gamma$;
then $PW_j\Vdash x_{\gamma+1}$ and \ref{lem:iterationedges} transports it
backwards.  At a nonzero limit the same $W_i$ realizes every earlier stage;
countable-supremum closure gives that full-prefix supremum.  This uses all
earlier stages, not just a cofinal subsequence of a nonmonotone orbit.
Finally beta-orient the displayed lambda term and apply
\ref{lem:realclosure}.
\end{proof}

\begin{lemma}[Zero, successor and a finite cap]\label{lem:primitives}
Closed native terms realize $0$, right successor and $\omega$.
For every $P:B\to B,U:B$, a closed term lies strictly above every $P^nU$.
\end{lemma}
\begin{proof}
$Z=\lambda z:o.z$ realizes zero.  The selector
$C_B(s,t)=(\lambda a:B.\lambda b:B.a)st$ lies strictly above $t$ by outer
Select; hence $\mathbf s=\lambda a:\sigma.C_\sigma(Z,a)$ realizes right
successor.  Put $\operatorname{Cap}(P,U)=R^B0(\lambda n:o.P)U$.
The carrier lemma gives the base $\operatorname{Cap}\succ_cU$ and iterates
the implication to $P^nU$.  Thus $\operatorname{Cap}(\mathbf s,Z)$ has rank
at least $\omega$ and realizes it.
\end{proof}

\subsection{Matching the published calibration}
The preceding primitives realize ordinal addition by iterating right successor,
multiplication by iterating $x\mapsto x+\alpha$ from zero, and
$\alpha\mapsto\omega^\alpha$ by iterating $x\mapsto x\omega$ from one.
For clarity this is an effective literal compilation, not an existence
claim about arbitrary semantic constants. With $\mathbf w=
\operatorname{Cap}(\mathbf s,Z)$ and $\mathbf1=\mathbf sZ$, the terms are
\[
 \mathbf{add}=\lambda a\,b.\It_\iota b\,\mathbf s\,a,\qquad
 \mathbf{mul}=\lambda a\,b.\It_\iota b\,
                    (\lambda x.\mathbf{add}\,x\,a)\,Z,
\]
\[
 \mathbf{pow}=\lambda a.\It_\iota a\,
                 (\lambda x.\mathbf{mul}\,x\,\mathbf w)\,\mathbf1.
\]
Their omitted binder types are all $\sigma$. The full-prefix iteration
theorem proves their stated ordinal meanings, including at limits.
They therefore realize
\[
 \operatorname{Fix}(f)(\zeta)=f^\omega(\zeta+1),\quad
 \operatorname{Next}=\operatorname{Fix}(\alpha\mapsto\omega^\alpha),
 \quad\epsilon=\operatorname{Next}(0).
\]
The literal realizer of $\operatorname{Fix}$ is
$\lambda f\,z.\It_\iota\mathbf w\,f\,(\mathbf s z)$;
apply it to $\mathbf{pow}$ for $\operatorname{Next}$ and then to $Z$
for $\epsilon$.
Let $A_0=\iota$, $A_{n+1}=A_n\to A_n$. For
$h:A_{l+1}$ and $h_j:A_j$, define $B_l:A_{l+2}$ by
\[
 B_lh\,h_l\cdots h_1
 =\operatorname{Fix}(\lambda\alpha^{\iota}.
                 h^\alpha h_l\cdots h_1\,0).
\]
For $l=0$ use $h^\alpha0$; for $l>0$ its first iteration is the instance
at $A_l$. Its realizer replaces $h^\alpha h_l$ by
$\It_{A_l}\alpha h h_l$ and applies the remaining arguments and $Z$;
at $l=0$ use $\It_\iota\alpha h Z$. Thus every $B_l$ has a finite native
term obtained by explicit lambda abstraction and application.
Define
\[
 d_0=\epsilon,\qquad d_1=\operatorname{Next}(\epsilon),\qquad
 d_{l+2}=B_lB_{l-1}\cdots B_0\operatorname{Next}\epsilon.
\]
These are the precise finite-type extractors of
\cite[Definitions 2.4--2.5, 2.10, 2.12, 2.14; Theorem 6.1 and Corollary 6.2]{simmons}.
Their helpfulness hypotheses are supplied by
\cite[Definition 5.1 and Theorem 5.10]{simmonshelpful}.
More precisely the semantic types are the full finite endotype hierarchy,
iteration takes the pointwise supremum of the \emph{entire} earlier prefix,
the fixed-point functional starts at $\zeta+1$, and the base pair is
$(\operatorname{Next},\varepsilon_0)$. These are exactly the hypotheses of
that calibration, including the terminal zero argument in $B_l$; a base
pair starting at $\omega$ is not substituted for it.

Let $a_0=0$, $a_{n+1}=\Omega^{a_n}$, with $\Omega=\omega_1$.
The published comparison identifies $d_n$ with $\psi_A(a_n)$ and
$\sup_na_n=\varepsilon_{\Omega+1}$. Here $\psi_A$ is the hull collapse
generated from the displayed base pair using addition, $\omega$-exponentiation
and available smaller collapses. The upper bound uses Pohlers's
Veblen-closed convention $\psi_B$. Simmons's comparison is cofinal agreement
of those conventions, not pointwise equality. Their monotonicity and the
cofinal tower sequence give equal final suprema. At the limiting argument,
each finite hull derivation uses finitely many collapse parameters, so is
already present below one earlier bound. Hence
\[
 \sup_nd_n=\sup_n\psi_A(a_n)
   =\psi_B(\varepsilon_{\Omega+1})=\BH.
\]
The cited calibration and helpfulness theorems supply this ordinal-analysis
step; the native carrier and iterator proofs above supply its new term
interpretation for the particular confined relation.

\begin{theorem}[Exact descending height]\label{thm:height}
Every individual $\succ_c$ rank is below $\BH$, and
\[
 \sup_s(\rk_{\succ_c}(s)+1)=\BH.
\]
An effective family of closed native terms of outer type $o\to o$ already
has cofinal ranks.
\end{theorem}
\begin{proof}
Replace every primitive in the finite expression $d_n$ by its native realizer
and apply the typed fundamental lemma.  This gives an effective finite term
$t_n:\sigma$ with $\rk(t_n)\ge d_n$.  There are only finitely many type
instances in each expression, and confinement has been checked for every
carrier, context and beta step.  Consequently
$\BH=\sup_nd_n\le\sup_n\rk(t_n)$, bounded above by the global height.
The individual upper theorem gives the reverse inequality.
\end{proof}

\section{Interpretation and limitations}
The construction has an exact ranked first-order RPO restriction and a literal
original-T subsystem.  It does not identify $d(n,a,b)$ with the native curried
$D^o n a b$.  The latter encoding has additional application-prefix structure,
and exact RPO agreement under it is a stronger question.  Nor does the height
theorem assert a total ordering or the order type of any linear extension.

The comparison algorithm is finite but can be extremely large: the domains
contain full monotone tuple tables.  No practical complexity bound, executable
implementation, or optimal rank notation is asserted.  Finally, this ordering
is not defined by comparing the numerical Wilken--Weiermann canonical vectors.
Its decidability leaves the residual arithmetic problem of
\cite{wwfoundations,wwexpressivity} separate.

\begin{thebibliography}{99}

\bibitem{ww}Gunnar Wilken and Andreas Weiermann.
Derivation lengths classification of G\"odel's T extending Howard's assignment.
\emph{Logical Methods in Computer Science} \textbf{8}(1:19) (2012), 1--44.
\href{https://doi.org/10.2168/LMCS-8(1:19)2012}
{doi:10.2168/LMCS-8(1:19)2012}.
\href{https://lmcs.episciences.org/1073}{Official article and full text}.
The 2012 revised article is the version used throughout.

\bibitem{kop}Cynthia Kop.
\emph{Higher Order Termination: Automatable Techniques for Proving
Termination of Higher-Order Term Rewriting Systems}.
PhD thesis, Vrije Universiteit Amsterdam, 2012.
IPA Dissertation Series 2012-14; ISBN 978-94-6203-164-7.
\href{https://www.cs.ru.nl/~cynthiakop/phdthesis.pdf}
{Author-hosted copy with corrections}.
Definition 5.19, printed pp.~117--118, is the recursive StarHORPO
definition used here.

\bibitem{simmons}Harold Simmons.
A comparison of two systems of ordinal notations.
\emph{Archive for Mathematical Logic} \textbf{43} (2004), 65--83.
\href{https://doi.org/10.1007/s00153-003-0177-z}
{doi:10.1007/s00153-003-0177-z}.
Definitions 2.4--2.5, 2.10, 2.12 and 2.14 specify the extractor
interface; Theorem 6.1 and Corollary 6.2 provide its calibration.

\bibitem{simmonshelpful}Harold Simmons.
Fruitful and helpful ordinal functions.
\emph{Archive for Mathematical Logic} \textbf{47} (2008), 677--709.
\href{https://doi.org/10.1007/s00153-008-0102-6}
{doi:10.1007/s00153-008-0102-6}.
Definition 5.1 and Theorem 5.10 supply the extractor and helpfulness
results used by the comparison calibration.

\bibitem{pohlers}Wolfram Pohlers.
\emph{Proof Theory: The First Step into Impredicativity}.
Universitext. Springer, Berlin and Heidelberg, 2009.
\href{https://doi.org/10.1007/978-3-540-69319-2}
{doi:10.1007/978-3-540-69319-2}.
The cited upper analysis uses Definition 9.1.2 and
Theorems 9.5.5--9.5.6; collapse conventions must be matched through the cited comparison theorems.

\bibitem{wwfoundations}\emph{The Wilken--Weiermann Comparison Problem: Foundations and a Guide to the Results}. Companion working manuscript, 2026.
\href{../05-wilken-weiermann-foundations/paper.pdf}{Included manuscript}.

\bibitem{wwexpressivity}\emph{Actual-Source Arithmetic and Undecidability Boundaries for Canonical Comparison}. Companion working manuscript, 2026.
\href{../09-canonical-source-expressivity/paper.pdf}{Included manuscript}.
\end{thebibliography}

\end{document}
